%%%%%%%%%%%%%%%%%%%%%%%%%%%%%%%%%
%
%       EXERCÍCIO
%
%%%%%%%%%%%%%%%%%%%%%%%%%%%%%%%%%

\ifdefstring{\atividade}{prova}{%
    \renewcommand{\valorquestao}{\ValorQ\ pontos}
}%

\begin{exercicioBanco}[\valorquestao]
Classifique cada uma das variáveis abaixo em qualitativa (nominal ou ordinal) ou quantitativa
(discreta ou contínua).
%
\begin{enumerate}[(a)]
    \item Ocorrência de hipertensão pré-natal em grávidas com mais de 35 anos (sim ou não).
    %
    \item Intenção de voto para presidente.
    %
    \item Perda de peso de maratonistas em uma corrida (em quilos).
    %
    \item Intensidade da perda de peso de maratonistas em uma corrida (leve, moderada, forte).
    %
    \item Grau de satisfação da população brasileira com relação ao trabalho de seu presidente (totalmente insatisfeito, insatisfeito, nem satisfeito nem insatisfeito, satisfeito, totalmente satisfeito).
\end{enumerate}
%
\end{exercicioBanco}

